% ---
% id: transmitter/decoy/rem:decoy-consumers
% kind: rem
% title: "Consumers of the emitted decoy declaration"
% summary: >-
%   The ensemble form of the consuming declarations reads the decoy declaration,
%   and exactly trusted-source mode, the rule against general constraint
%   entries, which refuses nothing it emits, and the emission's own refusals
%   stand in between. It licenses no restriction of a consumer's feasible set.
% prominence: standard
% status: accepted
% uses: [cor:emitted-decoy, req:decoy-emission]
% mentions: [transmitter/def:srcert, transmitter/def:source-spec, transmitter/rem:boundary, protocol-compiler/def:declaration, meat-core/def:declaration, protocol-compiler/def:eligible, reductions/def:chain, reductions/thm:two-sided, protocol-compiler/def:chaincond, req:entropy-cutoff, req:exact-decoy-blocks]
% assumes: []
% refs: []
% provenance: newly-derived
% lean:
% ---

\begin{remark}\label{rem:decoy-consumers}
By Corollary~\ref{cor:emitted-decoy}\,(i) and (ii) the decoy declaration is a
source specification in trusted-source mode given by an ensemble presentation,
and its source segment meets the admission conditions of the transmitter inside
the artefact the emission outputs (Requirement~\ref{req:decoy-emission}). Item
(a) says which consumer reads it; what stands between the decoy declaration and
a consumer of that form is exactly (b) to (d), each a named refusal.
\begin{enumerate}[label=(\alph*),leftmargin=2.2em]
\item \emph{The consumer's form.} The protocol compiler's declaration
(\xref{protocol-compiler/def:declaration}) and the MEAT core's
(\xref{meat-core/def:declaration}) each have an ensemble form, which carries
$\{N_x\}$ in the place of Alice's POVM, and
\xref{protocol-compiler/def:eligible} admits an ensemble presentation beside a
declaration of that form, so the decoy declaration has a consumer. A consumer
whose declaration carries $\{M_x\}$ refuses it (Definition~\ref{def:srcert}),
which is a refusal of that consumer's form and not a condition on the decoy
declaration.
\item \emph{Trusted-source mode only.} A non-empty source segment is admitted in
trusted-source mode alone, by the admission rule of \xref{reductions/def:chain},
and the two-sided reduction theorem (\xref{reductions/thm:two-sided}) is stated
for that mode alone; an ensemble form is refused in
\tref{untrusted-source mode}, the
mode without a marginal or a certificate (Definition~\ref{def:source-spec}), by
the consuming module.
\item \emph{Scalar constraint entries only.} Beside a non-empty source segment
the protocol compiler requires every entry of a consumer's constraint list
$\CList$ to have a real scalar as its right-hand side, and refuses an entry
declared general (the licensed-only rule of
\xref{protocol-compiler/def:chaincond}). The decoy declaration carries no
constraint entry (Corollary~\ref{cor:emitted-decoy}\,(ii)), so the rule refuses
nothing it emits.
\item \emph{The emission's own refusals.} The requirements of this submodule
refuse a declaration before any consumer reads it; the bound
$N_{\mathrm{ent}}\le2$ on the entropy cutoff
(Requirement~\ref{req:entropy-cutoff}) and the exactness of the formed entries
(Requirement~\ref{req:exact-decoy-blocks}) are the two of them a declarer of a
physically interesting source meets first.
\end{enumerate}
One further point is not a refusal but the absence of a licence: the emission
licenses no restriction of any consumer's feasible set. The states reached from
its certificate are block diagonal, but a positive operator constrained only to
have the marginal $\tau_A$ need not be (Corollary~\ref{cor:emitted-decoy}\,(iii)),
and a consumer that restricts its own optimisation does so under a theorem of
its own module (Remark~\ref{rem:boundary}).
\end{remark}
